\begin{table*}
    \scriptsize
    \centering
    \extrarowheight=\aboverulesep
    \addtolength{\extrarowheight}{\belowrulesep}
    \aboverulesep=0pt
    \belowrulesep=0pt
    \setlength{\tabcolsep}{0.45em}
        \begin{tabular}{>{\centering\arraybackslash\columncolor{Gray}}m{1.0cm} x{3.5cm} y{3.5cm} x{4.7cm} y{2.3cm}}
        \toprule
        \textbf{Success} & \textbf{Utterance 1} & \textbf{Utterance 2} & \textbf{Concatenation} & \textbf{Implication} \\
        \midrule
         X & weather for frankfort (\texttt{GetWeather}) & 3 out of 6 for the last album (\texttt{RateBook}) & frankfort's weather gets a 3 out of 6 rating for the last album & Intent (\texttt{RateBook}) removed \\
         & show me the fares for a first class ticket from baltimore to dallas (\texttt{atis\_airfare}) & i would like the time of all flights from san francisco to pittsburgh on sunday (\texttt{atis\_flight\_time}) & show me the fares for a first class ticket from baltimore to dallas and i would like to know all flights from san francisco to pittsburgh on sunday & Intent changed (from \texttt{atis\_flight\_time} to \texttt{atis\_flight}) \\
        & what is mci (\texttt{atis\_abbreviation}) & list la (\texttt{atis\_city}) & combine the sentences: "what is mci?" and "list la". & Failed to merge \\
        \midrule
        O & play the playlist funtime activity 
        (\texttt{AddToPlaylist}) & add adele onto funtime activity playlist
        (\texttt{PlayMusic}) & play the playlist, adding adele to the funtime activity playlist & Gerund phrase \\
        & how do i freeze my bank account	
        (\texttt{account\_blocked}) & do you know why my bank account is frozen (\texttt{freeze\_account}) & how do i freeze my bank account, and do you know why it is frozen & Coreference \\
        \bottomrule
        \end{tabular}
    \caption{Failed and successful results of the Generative Approach and their implications.}
    \label{table:generative}
\end{table*}

\subsection{Concatenation}
\label{subsec:concatenation}
MID datasets are typically created by merging two utterances using connectives. 
However, this fails to encompass the full range of ways people express multiple intents, as they often employ varied connectives or omit them entirely.
To assemble multi-intent utterances with nuanced patterns and to improve upon the rule-based approach suggested by MixX, we view concatenation from two distinct aspects, as in Figure \ref{fig:concat}: 
\textbf{the complexity of concatenation} (explicit or implicit) and \textbf{the methodology of performing concatenation} (whether conducted manually or through tools such as ChatGPT.)

\paragraph{The complexity aspect}
We present two merging methods based on their level of complexity.

\begin{enumerate}
    \item \textbf{Explicit Concatenation:}    
    Conjunctions are explicitly used to concatenate two or more utterances.\footnote{Specifically: \texttt{`and'}, \texttt{`and then'}, \texttt{`and also'}, \texttt{`,(comma)'}, \texttt{`;(semi-colon)'}, \texttt{`or'}, \texttt{`before'}, \texttt{`after'}, \texttt{`additionally'}, \texttt{`finally'}.\label{fn:conjunctions}}
    We refer to the use of the four connectives as outlined in \citetlanguageresource{qin-etal-2020-agif}---\texttt{`and'}, \texttt{`and then'}, \texttt{`and also'}, \texttt{`,(comma)'}---as the \textbf{AND variants} setting.    
    If other conjunctions are employed, we denote it as \textbf{various conjunctions} (see Figure \ref{fig:concat}).
    \item \textbf{Implicit Concatenation:} This approach pursues a seamless blend of utterances, minimizing the apparent usage of conjunctions.
    \citet{meng-etal-2022-dialogusr_} discovered that 62.5\% of the follow-up queries in the dialogues they gathered were either incomplete or incorrectly formulated.
    This inspires us to consider the four following implicit merging patterns.
    \textbf{Inherent ambiguity (conjunction removal)} refers to cases where conjunctions are simply removed from their original positions.
    Despite its simplicity, it effectively reflects the intuition that speakers tend to favor shorter utterances.
    In line with the same philosophy, we also consider \textbf{omissions} and \textbf{coreferences}, where redundant expressions are either eliminated or substituted with pronouns.
    Lastly, we employ \textbf{gerund phrases} (the \texttt{-ing} form of verbs) which are useful for emphasizing concurrency.
    For a clearer understanding of readers, we provide examples of each case in Table \ref{table:concat_and_metrics}.
\end{enumerate}

\paragraph{The methodology aspect}
The remaining issue is how we implement the phenomena we have specified.
While we permit minor adjustments to source sentences during merging, like omissions and coreferences, we aim to retain the sentences' original structure to the greatest extent possible.
In alignment with our mission, we propose two implementation strategies: manual rule-based heuristics and the utilization of ChatGPT.

\begin{enumerate}
    \item \textbf{Na\"ive Approach:}
    It follows the original practices proposed in \citetlanguageresource{qin-etal-2020-agif}.
    It is Explicit Concatenation with the AND variants setting. 
    \item \textbf{Manual Approach:} It expands the Na\"ive Approach with further rule-based techniques. 
    To be specific, it utilizes a broader range of conjunctions (see Footnote \ref{fn:conjunctions}), occasionally skips connectives, and, when feasible, transforms one of the utterances into a gerund phrase.
    \item \textbf{Generative Approach}: Facing the non-trivial challenge of implicit concatenations, we attempt to circumvent this issue with the aid of large language models, especially ChatGPT.
    We instruct the tool to merge utterances while preserving their original intents and structures as much as possible, and we specifically encourage it to avoid using conjunctions such as \texttt{'and'}.
    Figure \ref{figure:prompt_chatgpt_concat} shows the prompt used in the process. 
    We provide three few-shot samples each of both successful and unsuccessful merges.
    More detailed information about the prompt is available in Appendix \ref{subsec:prompt_concat}.
    \label{table:metric_samples}
\end{enumerate}

\begin{figure}[t]
    \setlength{\fboxsep}{7pt} 
    \noindent\fbox{%
        \begin{minipage}{\dimexpr\columnwidth-2\fboxsep-2\fboxrule\relax}
        
            \fontsize{9pt}{11pt}\selectfont 
            \setlength{\parskip}{5pt} 
            \textbf{Instructions:} \textsc{You are a native English speaker. Combine the following sentences as one single sentence naturally.} \\
            \vspace{-5pt}

            \textcolor{commentcolor}{\# $N$ iterations} \\
            \vspace{-5pt}
            \textbf{Example $N$: } [($\mathrm{utt_1}$, $\mathrm{intent_1}$), ..., ($\mathrm{utt_k}$, $\mathrm{intent_k}$)]
            \begin{itemize}
                \item \textbf{Good Answer: } \textsc{[sample answer]} \\\vspace{-3pt}\textcolor{commentcolor}{\# An ideal result of concatenating example utterances}
                \item \textbf{Bad Answer: } \textsc{[sample answer]} \\\vspace{-3pt}\textcolor{commentcolor}{\# An unwanted or incorrect result of concatenating example utterances}
            \end{itemize}

            \textbf{Query: } \textsc{[query]} \\\vspace{-5pt}
            
            \textbf{Answer: } \textsc{[answer]}
        \end{minipage}%
    }
    \caption{
    Prompt design for the Generative Approach. 
    The demonstration section showcases $N$ ($=$ 3)  examples of combining $k$ ($=$ 2 or 3) utterances, featuring both a successful and a failed case. 
    The \textsc{[query]} lists the sentences to be merged.
    ChatGPT performs the merging process by filling in the \textsc{[answer]} part. 
    Blue comments are for illustrative purposes only and are not part of the actual prompt.
    }
    \label{figure:prompt_chatgpt_concat}
\end{figure}

In our preliminary experiments, we discovered that ChatGPT often fails to combine two given sentences with the expected level of creativity and naturalness.
Consequently, contrary to our expectations, the initial results from ChatGPT closely reproduced those from the Manual Approach, which heavily relies on directly using conjunctions.
For reference, Table \ref{table:generative} displays a few examples of failed and successful outputs from ChatGPT.

Nevertheless, we noted some instances where ChatGPT excels in producing high-quality utterances (e.g., the bottom part of Table \ref{table:generative}) that could not be simply achieved through rule-based heuristics.
We highlight that manually designing such examples requires an exceptional level of expertise and effort, and they are not readily achievable even by human annotators.
Therefore, in the subsequent section, we investigate techniques in terms of utterance selection to better condition ChatGPT for creating more reliable data samples.
